% Provenance (never printed): NSI terminale "Structures de donnees: graphes (listes d'adjacence, matrice);
% algorithmique: parcours en profondeur et en largeur, detection de cycle, chemin" and the Dijkstra algorithm met in
% the routing chapter. See sources/book2/BRIEFS.md.
\chapter{Graph Traversals and Shortest Paths}\label{ch:g12:graph-algorithms}

A classic way out of a maze is to put one hand on a wall and never let go. It works in every maze whose exit is on the
outer wall, and it needs no memory at all. But it walks into every dead end on its side, and comes back out of each:
the way it finds is rarely the shortest.

This chapter writes the two systematic ways of exploring a \omterm{def:g10:graphs-first-look:graph}{graph}, depth first with a stack and breadth first with a
queue, and uses them to decide what can be reached, to find \omterm{def:g12:graph-algorithms:cycle}{cycles} and to find shortest paths counted in edges. Then
the edges get lengths, and \omterm{met:g12:graph-algorithms:dijkstra}{Dijkstra's algorithm} finds the shortest routes across the town.

\begin{recall}
A \omterm{def:g10:graphs-first-look:graph}{graph} is made of vertices and edges, possibly directed and weighted; a path, its length, the distance between two
vertices, and the components of a \omterm{def:g10:graphs-first-look:graph}{graph} were defined in chapter \ref{ch:g10:graphs-first-look}, where the town's road
\omterm{def:g10:graphs-first-look:graph}{graph} was first loaded. A stack gives back its newest element first, a queue its oldest (chapter
\ref{ch:g12:lists-stacks-queues}). The adjacency matrix of a \omterm{def:g10:graphs-first-look:graph}{graph} (the mathematics book).
\end{recall}

\section{Representations}\label{sec:g12:graph-algorithms:representations}

\begin{definition}[Adjacency list, successor, predecessor]\label{def:g12:graph-algorithms:adjacency-list}
The \emph{adjacency lists}\index{adjacency list} of a \omterm{def:g10:graphs-first-look:graph}{graph} give, for each vertex, the list of its neighbours; in
Python, a dictionary from each vertex to that list. In a \omterm{def:g10:graphs-first-look:directed}{directed graph}, the vertices at the end of the edges leaving
$v$ are its \emph{successors}\index{successor}, and those at the start of the edges arriving at $v$ its
\emph{predecessors}\index{predecessor}; the lists usually give the successors.
\end{definition}

The adjacency matrix of a \omterm{def:g10:graphs-first-look:graph}{graph} of $n$ vertices is an $n \times n$ table of $0$ and $1$, with a $1$ in row $v$ and
column $w$ when there is an edge from $v$ to $w$.

\omcode{code/grade-12/05-graph-algorithms/listings/representations.py}{1}{29}{The two representations, converted both
ways, and the \omterm{def:g12:graph-algorithms:adjacency-list}{predecessors} of a \omterm{def:g10:graphs-first-look:directed}{directed graph}.}\label{lst:g12:graph-algorithms:representations}

\omoutput{code/grade-12/05-graph-algorithms/outputs/representations.txt}{Each row of the matrix holds the same
information as one list.}\label{out:g12:graph-algorithms:representations}

\begin{omfigure}
\begin{tikzpicture}[font=\footnotesize, v/.style={draw, circle, minimum size=6mm, inner sep=0pt, fill=omDef!10}]
  \node[v] (A) at (0,1) {A}; \node[v] (B) at (1.2,2) {B}; \node[v] (C) at (1.2,0) {C};
  \node[v] (D) at (2.6,1.4) {D}; \node[v] (E) at (2.6,-0.4) {E}; \node[v] (F) at (3.8,0.6) {F};
  \foreach \a/\b in {A/B, A/C, B/D, C/D, C/E, D/F, E/F} \draw (\a) -- (\b);
  \node[align=left, right] at (4.6,0.8) {\code{A: B C}\\ \code{B: A D}\\ \code{C: A D E}\\ \code{D: B C F}\\
    \code{E: C F}\\ \code{F: D E}};
  \matrix[matrix of math nodes, right, nodes={minimum width=4mm, font=\scriptsize}, left delimiter={(},
      right delimiter={)}] at (7.6,0.8) {
    0 & 1 & 1 & 0 & 0 & 0 \\ 1 & 0 & 0 & 1 & 0 & 0 \\ 1 & 0 & 0 & 1 & 1 & 0 \\
    0 & 1 & 1 & 0 & 0 & 1 \\ 0 & 0 & 1 & 0 & 0 & 1 \\ 0 & 0 & 0 & 1 & 1 & 0 \\};
\end{tikzpicture}

{\small The \omterm{def:g10:graphs-first-look:graph}{graph} \code{G} of the listings, its \omterm{def:g12:graph-algorithms:adjacency-list}{adjacency lists} and its adjacency matrix (rows and columns in the order
A to F). The matrix is symmetric because the \omterm{def:g10:graphs-first-look:graph}{graph} is undirected.}
\end{omfigure}

The two cost differently. The matrix always takes $n^2$ cells, and answers ``is there an edge from $v$ to $w$?'' in one
step. The lists take one entry per edge end, and list the neighbours of $v$ without looking at the others. For a town
of $30$ places and $53$ roads the lists are far smaller; most \omterm{def:g10:graphs-first-look:graph}{graphs} met in practice are like it, with few edges per
vertex, and are stored as lists.

\section{Depth-first search}\label{sec:g12:graph-algorithms:dfs}

\begin{method}[Depth-first search]\label{def:g12:graph-algorithms:dfs}
A \emph{depth-first search}\index{depth-first search} explores a \omterm{def:g10:graphs-first-look:graph}{graph} from a starting vertex by always going on from
the most recently found vertex: visit it, mark it as visited, then search from each of its unvisited neighbours in
turn, before coming back. A stack holds the vertices still to explore; the recursion's \omterm{def:g12:recursion:call-stack}{call stack} can play that role.
\end{method}

\omcode{code/grade-12/05-graph-algorithms/listings/dfs.py}{1}{33}{Depth-first search with a stack and with
recursion.}\label{lst:g12:graph-algorithms:dfs}

\omoutput{code/grade-12/05-graph-algorithms/outputs/dfs.txt}{Both versions give the same order; everything is
reachable from E.}\label{out:g12:graph-algorithms:dfs}

\begin{omfigure}
\begin{tikzpicture}[font=\footnotesize, v/.style={draw, circle, minimum size=6mm, inner sep=0pt, fill=omDef!10}]
  \node[v] (A) at (0,1) {A}; \node[v] (B) at (1.6,2.2) {B}; \node[v] (C) at (1.6,-0.2) {C};
  \node[v] (D) at (3.4,1.6) {D}; \node[v] (E) at (3.4,-0.6) {E}; \node[v] (F) at (5.0,0.6) {F};
  \foreach \a/\b in {A/B, B/D, D/C, C/E, E/F} \draw[very thick, omDef] (\a) -- (\b);
  \foreach \a/\b in {A/C, D/F} \draw[dashed, black!50] (\a) -- (\b);
  \foreach \n/\k/\x/\y in {A/1/-0.45/0, B/2/0/0.45, D/3/0/0.45, C/4/0/-0.45, E/5/0/-0.45, F/6/0.45/0}
    \node[omThm, font=\scriptsize\bfseries] at ($(\n)+(\x,\y)$) {\k};
\end{tikzpicture}

{\small The depth-first search of \code{G} from A: the numbers give the order of visit. Bold edges are those used to
reach a new vertex; they form a tree. The dashed edges lead to vertices already visited.}
\end{omfigure}

A search from $v$ visits exactly the vertices reachable from $v$; each vertex is visited once, and each edge looked at
once from each end, so the cost is proportional to the number of vertices plus the number of edges.

\begin{definition}[Cycle]\label{def:g12:graph-algorithms:cycle}
A \emph{cycle}\index{cycle} is a path of at least one edge that comes back to its first vertex without using any edge
twice; in an undirected \omterm{def:g10:graphs-first-look:graph}{graph}, it has at least three vertices.
\end{definition}

In an undirected \omterm{def:g10:graphs-first-look:graph}{graph}, a depth-first search that meets an already visited vertex, other than the one it came from,
has found a \omterm{def:g12:graph-algorithms:cycle}{cycle}. In a \omterm{def:g10:graphs-first-look:directed}{directed graph} that is not enough: two paths may join without any \omterm{def:g12:graph-algorithms:cycle}{cycle}. The search then gives
each vertex one of three colours: white, not yet seen; grey, being explored, its search not finished; black,
finished. An edge to a grey vertex goes back to a vertex whose search is still running, so there is a path from it
to the present vertex and back: a \omterm{def:g12:graph-algorithms:cycle}{cycle}.

\omcode{code/grade-12/05-graph-algorithms/listings/cycle.py}{1}{22}{\omterm{def:g12:graph-algorithms:cycle}{Cycles} in an undirected
\omterm{def:g10:graphs-first-look:graph}{graph}.}\label{lst:g12:graph-algorithms:cycle-u}

\omcode{code/grade-12/05-graph-algorithms/listings/cycle.py}{23}{46}{\omterm{def:g12:graph-algorithms:cycle}{Cycles} in a \omterm{def:g10:graphs-first-look:directed}{directed graph}, with three colours;
course prerequisites.}\label{lst:g12:graph-algorithms:cycle-d}

\omoutput{code/grade-12/05-graph-algorithms/outputs/cycle.txt}{A course that would require itself is a
\omterm{def:g12:graph-algorithms:cycle}{cycle}.}\label{out:g12:graph-algorithms:cycle}

\section{Breadth-first search}\label{sec:g12:graph-algorithms:bfs}

\begin{method}[Breadth-first search]\label{def:g12:graph-algorithms:bfs}
A \emph{breadth-first search}\index{breadth-first search} from a vertex $s$ visits $s$, then all its neighbours, then
all their unvisited neighbours, and so on, layer by layer. A queue holds the vertices found but not yet explored. Each
vertex records its distance from $s$ in edges and its parent, the vertex from which it was found.
\end{method}

\omcode{code/grade-12/05-graph-algorithms/listings/bfs.py}{1}{32}{Breadth-first search: distances and
parents.}\label{lst:g12:graph-algorithms:bfs}

\omoutput{code/grade-12/05-graph-algorithms/outputs/bfs.txt}{Following the parents back from F gives a shortest
path.}\label{out:g12:graph-algorithms:bfs}

The queue gives back the vertices in order of distance: all those at distance $d$ before any at distance $d + 1$. So
the first time a vertex is found, it is found by a shortest path (exercise \ref{exo:g12:graph-algorithms:12}). The
parents form a tree, the breadth-first search tree, and the path from $s$ to any vertex in that tree is a shortest one.

A maze drawn on squared paper is a \omterm{def:g10:graphs-first-look:graph}{graph}: each square is a vertex, and two squares are joined when no wall separates
them.

\begin{omfigure}
\begin{tikzpicture}[x=10mm, y=-10mm, font=\footnotesize]
  \foreach \r/\c/\d in {0/0/0, 0/1/1, 0/2/2, 0/3/5, 1/0/1, 1/1/4, 1/2/3, 1/3/4, 2/0/2, 2/1/5, 2/2/6, 2/3/7} {
    \pgfmathsetmacro{\shade}{45 - 5*\d}
    \fill[omDef!\shade!white] (\c,\r) rectangle (\c+1,\r+1);
    \node at (\c+0.5,\r+0.5) {$\d$};
  }
  \draw[thick] (0,0) rectangle (4,3);
  \draw[very thick] (1,1) -- (2,1); \draw[very thick] (1,1) -- (1,2); \draw[very thick] (2,2) -- (3,2);
  \draw[very thick] (3,0) -- (3,1); \draw[very thick] (3,2) -- (4,2); \draw[very thick] (1,2) -- (1,3);
\end{tikzpicture}

{\small A small maze and the distances found by breadth-first search from the top-left square. Each number is one
more than a neighbour's that can be reached without crossing a wall; the darker the square, the closer.}
\end{omfigure}

The program below builds a maze of $8 \times 14$ squares by carving passages at random from a fixed seed, then knocks
down twelve more walls so that some places can be reached in several ways. It finds the shortest way from the
top-left square to the bottom-right one, and compares it with the hand on the wall.

\omcode{code/grade-12/05-graph-algorithms/listings/maze.py}{1}{35}{A maze as a grid \omterm{def:g10:graphs-first-look:graph}{graph}, carved by a depth-first
search.}\label{lst:g12:graph-algorithms:maze}

\omcode{code/grade-12/05-graph-algorithms/listings/maze.py}{48}{70}{The right hand on the wall, and the
comparison.}\label{lst:g12:graph-algorithms:maze-run}

\omoutput{code/grade-12/05-graph-algorithms/outputs/maze.txt}{The stars mark the shortest exit. Lines 36 to 47 of the
file, not printed, draw the maze as text.}\label{out:g12:graph-algorithms:maze}

\section{Dijkstra's algorithm}\label{sec:g12:graph-algorithms:dijkstra}

When edges have lengths, the path with fewest edges is not always the shortest: breadth-first search no longer
answers. The \omterm{def:g10:graphs-first-look:path}{length of a path} is now the sum of the weights of its edges, all assumed positive or zero.

\begin{method}[Dijkstra's algorithm]\label{met:g12:graph-algorithms:dijkstra}
\emph{Dijkstra's algorithm}\index{Dijkstra's algorithm} computes the distances from a source $s$ in a \omterm{def:g10:graphs-first-look:graph}{graph} with
weights $\ge 0$. It keeps a tentative distance for each vertex found, starting with $0$ for $s$, and a set of
finished vertices, at first empty. At each step:
\begin{enumerate}
  \item among the vertices found and not finished, pick one, $u$, with the smallest tentative distance, and finish it;
  \item for each edge from $u$ to an unfinished $v$ of weight $w$, if $d(u) + w$ is smaller than the tentative
    distance of $v$, or $v$ has none, set it to $d(u) + w$, and record $u$ as the parent of $v$.
\end{enumerate}
It stops when every vertex found is finished.
\end{method}

\begin{history}[Dijkstra, 1959]
Edsger W.\ Dijkstra published the method in 1959, in a three-page note titled ``A note on two problems in connexion
with \omterm{def:g10:graphs-first-look:graph}{graphs}'', in the first volume of the journal \emph{Numerische Mathematik}.
% ledger: g12:dijkstra-1959
\end{history}

\omcode{code/grade-12/05-graph-algorithms/listings/dijkstra.py}{1}{30}{\omterm{met:g12:graph-algorithms:dijkstra}{Dijkstra's algorithm}, with the minimum found by
a scan.}\label{lst:g12:graph-algorithms:dijkstra}

\omoutput{code/grade-12/05-graph-algorithms/outputs/dijkstra.txt}{The tentative distances after each vertex is
finished.}\label{out:g12:graph-algorithms:dijkstra}

\begin{omfigure}
\begin{tikzpicture}[font=\footnotesize, v/.style={draw, circle, minimum size=6mm, inner sep=0pt, fill=omMeth!15}]
  \node[v] (S) at (0,1) {S}; \node[v] (A) at (2,2) {A}; \node[v] (B) at (2,0) {B};
  \node[v] (C) at (4.4,2) {C}; \node[v] (D) at (4.4,0) {D}; \node[v] (T) at (6.6,1) {T};
  \foreach \a/\b/\w/\t in {S/A/4/0, S/B/1/1, A/B/2/1, A/C/5/1, B/D/8/1, C/D/1/0, C/T/3/1, D/T/6/0} {
    \ifnum\t=1 \draw[very thick, omDef] (\a) -- node[fill=white, inner sep=1pt, text=black] {$\w$} (\b);
    \else \draw[black!50] (\a) -- node[fill=white, inner sep=1pt, text=black] {$\w$} (\b); \fi
  }
  \foreach \n/\d/\k/\x/\y in {S/0/1/-1.25/0, B/1/2/0/-0.5, A/3/3/0/0.5, C/8/4/0/0.5, D/9/5/0/-0.5, T/11/6/1.35/0}
    \node[omThm, font=\scriptsize] at ($(\n)+(\x,\y)$) {$d = \d$, \#\k};
\end{tikzpicture}

{\small The \omterm{def:g10:graphs-first-look:graph}{graph} \code{W}: each vertex with its final distance from S and the step at which it was finished. The bold
edges join each vertex to its parent; they form the tree of shortest routes. D is at $9$ through B and through C alike; the program keeps B, the first
parent found, since it changes a parent only for a strictly shorter route. B is finished before A although S--A is
a direct edge: through B, A is at $3$, not $4$.}
\end{omfigure}

\begin{theorem}[The finished distances are final]\label{thm:g12:graph-algorithms:dijkstra}
With weights $\ge 0$, when \omterm{met:g12:graph-algorithms:dijkstra}{Dijkstra's algorithm} finishes a vertex, its tentative distance is its distance from $s$.
\end{theorem}

The idea of the proof (exercise \ref{exo:g12:graph-algorithms:13}): take a shortest path from $s$ to the vertex $u$
being finished, and the first vertex $y$ on it that is not yet finished. The vertex before $y$ is
finished, and gave $y$ a tentative distance at most the length of the path up to $y$; that $d(y)$ is at least $d(u)$,
since $u$ was picked as the smallest; and the rest of the path has weights $\ge 0$. So
no path is shorter than $d(u)$. This version scans the unfinished vertices at each step: with $n$ vertices, $n$ steps
of up to $n$ comparisons, a cost in $n^2$. A better structure for picking the minimum, met in a later volume, makes
it much faster on large \omterm{def:g10:graphs-first-look:graph}{graphs}.

Negative weights break the proof, and the algorithm:

\omcode{code/grade-12/05-graph-algorithms/listings/negative_weight.py}{1}{22}{A negative weight, and every path tried
by brute force.}\label{lst:g12:graph-algorithms:negative}

\omoutput{code/grade-12/05-graph-algorithms/outputs/negative_weight.txt}{A finished at $1$, before the edge of weight
$-3$ was seen.}\label{out:g12:graph-algorithms:negative}

\section{On the town}\label{sec:g12:graph-algorithms:town}

The town's road \omterm{def:g10:graphs-first-look:graph}{graph} has $30$ places and $53$ roads, one of them one-way: road $1$ goes from the Town Square to the
Library only. \omterm{met:g12:graph-algorithms:dijkstra}{Dijkstra's algorithm} from the Station gives the length of the shortest drive to every place.

\omcode{code/grade-12/05-graph-algorithms/listings/dijkstra_town.py}{1}{21}{Shortest drives from the
Station.}\label{lst:g12:graph-algorithms:town}

\omoutput{code/grade-12/05-graph-algorithms/outputs/dijkstra_town.txt}{The one-way road makes the return from the
Library almost twice as long.}\label{out:g12:graph-algorithms:town}

\begin{omfigure}
\begin{tikzpicture}
\begin{axis}[omrecord, width=10cm, height=8.6cm, xlabel={$x$ (km)}, ylabel={$y$ (km)}, xmin=-0.2, xmax=6, ymin=-0.2,
    ymax=5.1, axis equal image, grid=none]
  \addplot[black!25, line width=0.6pt] table[x=x, y=y]{code/grade-12/05-graph-algorithms/figures/town_roads.dat};
  \addplot[omDef, line width=1.4pt] table[x=x, y=y]{code/grade-12/05-graph-algorithms/figures/town_tree.dat};
  \addplot[only marks, mark=*, mark size=1.6pt, omThm, nodes near coords, point meta=explicit symbolic,
    nodes near coords style={font=\tiny, black!70, anchor=south west, inner sep=1pt}]
    table[x=x, y=y, meta=d]{code/grade-12/05-graph-algorithms/figures/town_places.dat};
  \node[font=\scriptsize, anchor=north] at (axis cs:2.301,2.80) {Station};
\end{axis}
\end{tikzpicture}

{\small The town's roads (grey) and the tree of shortest routes from the Station (blue), drawn from the coordinates of
the places by \code{figures/town_tree.py}; each place is labelled with its distance from the Station in kilometres.}
\end{omfigure}

The functions above take a \omterm{def:g10:graphs-first-look:graph}{graph} as a dictionary. A class can hide that dictionary behind methods; if it offers
\code{g[v]} and a loop over the vertices, through two \omterm{def:g12:classes-and-objects:special-method}{special methods}, the same \code{dijkstra} works on it unchanged.

\omcode{code/grade-12/05-graph-algorithms/listings/graph_class.py}{1}{28}{A \omterm{def:g10:graphs-first-look:graph}{graph} class.}%
\label{lst:g12:graph-algorithms:graph-class}

\omcode{code/grade-12/05-graph-algorithms/listings/use_graph.py}{1}{11}{The \omterm{def:g10:graphs-first-look:graph}{graph} \code{W} rebuilt with the class, plus
a one-way edge.}\label{lst:g12:graph-algorithms:use-graph}

\omoutput{code/grade-12/05-graph-algorithms/outputs/use_graph.txt}{The same distances as before, and X reached through
its one-way edge.}\label{out:g12:graph-algorithms:use-graph}

\section{Exercises}\label{sec:g12:graph-algorithms:exercises}

\begin{exercise}[$\star$]\label{exo:g12:graph-algorithms:1}
Give the order of the recursive depth-first search of \code{G} from D.
\end{exercise}

\begin{exercise}[$\star$]\label{exo:g12:graph-algorithms:2}
Give the order of the breadth-first search of \code{G} from D, layer by layer.
\end{exercise}

\begin{exercise}[$\star$]\label{exo:g12:graph-algorithms:3}
Four vertices $1$, $2$, $3$, $4$ form a square, $1$--$2$--$3$--$4$--$1$, with one diagonal $1$--$3$. Write its
adjacency matrix and its \omterm{def:g12:graph-algorithms:adjacency-list}{adjacency lists}.
\end{exercise}

\begin{exercise}[$\star$]\label{exo:g12:graph-algorithms:4}
Give the distances in edges from E to every vertex of \code{G}.
\end{exercise}

\begin{exercise}[$\star$]\label{exo:g12:graph-algorithms:5}
Write \code{count_edges(adj)}, the number of edges of an undirected \omterm{def:g10:graphs-first-look:graph}{graph} given by \omterm{def:g12:graph-algorithms:adjacency-list}{adjacency lists}.
\end{exercise}

\begin{exercise}[$\star\star$]\label{exo:g12:graph-algorithms:6}
A student's search marks a vertex when it enters it and unmarks it when it returns, ``to keep the marks clean''. On a
grid of $4 \times 4$ squares, it makes $2111$ calls from a corner. How many does the depth-first search of the chapter
make, and why does the student's version make so many?
\end{exercise}

\begin{exercise}[$\star\star$]\label{exo:g12:graph-algorithms:7}
Write \code{has_path(g, a, b)}, which stops as soon as \code{b} is found.
\end{exercise}

\begin{exercise}[$\star\star$]\label{exo:g12:graph-algorithms:8}
Run \omterm{met:g12:graph-algorithms:dijkstra}{Dijkstra's algorithm} by hand on \code{W} from A: the vertex finished at each step and the tentative distances.
\end{exercise}

\begin{exercise}[$\star\star$]\label{exo:g12:graph-algorithms:9}
A version of \omterm{met:g12:graph-algorithms:dijkstra}{Dijkstra's algorithm} finishes, at each step, the vertex found last instead of the closest. On \code{W}
from S it gives $15$ for T. Find where it goes wrong.
\end{exercise}

\begin{exercise}[$\star\star$]\label{exo:g12:graph-algorithms:10}
The town has $30$ places and $53$ roads, one of them one-way. How many entries do its \omterm{def:g12:graph-algorithms:adjacency-list}{adjacency lists} hold, and how
many cells its adjacency matrix?
\end{exercise}

\begin{exercise}[$\star\star$]\label{exo:g12:graph-algorithms:11}
Write \code{bfs_layers(g, s)}, which returns the list of layers: \code{[[s], [vertices at distance 1], ...]}.
\end{exercise}

\begin{exercise}[$\star\star\star$]\label{exo:g12:graph-algorithms:12}
Prove that breadth-first search gives every reachable vertex its distance in edges from $s$.
\end{exercise}

\begin{exercise}[$\star\star\star$]\label{exo:g12:graph-algorithms:13}
Prove the theorem on \omterm{met:g12:graph-algorithms:dijkstra}{Dijkstra's algorithm}.
\end{exercise}

\begin{exercise}[$\star\star\star$]\label{exo:g12:graph-algorithms:14}
Write \code{topological_order(g)}, which orders the vertices of a \omterm{def:g10:graphs-first-look:directed}{directed graph} so that every edge goes from an
earlier vertex to a later one, and raises \code{ValueError} if the \omterm{def:g10:graphs-first-look:graph}{graph} has a \omterm{def:g12:graph-algorithms:cycle}{cycle}.
\end{exercise}

\begin{exercise}[$\star\star\star$]\label{exo:g12:graph-algorithms:15}
To remove negative weights, a student adds the same number to every weight, then runs \omterm{met:g12:graph-algorithms:dijkstra}{Dijkstra's algorithm}. Show with
the \omterm{def:g10:graphs-first-look:graph}{graph} of \code{negative_weight.py} that the path found is not the shortest in the original \omterm{def:g10:graphs-first-look:graph}{graph}.
\end{exercise}

\section{Problem: The Shortest Exit}\label{sec:g12:graph-algorithms:problem}

\begin{problem}[{Weekend problem --- the shortest exit}]\label{pb:g12:graph-algorithms:1}
The maze of \code{maze.py} has $8 \times 14$ squares, an entrance at the top-left square and an exit at the
bottom-right one. In its second version, $30$ of its passages, chosen at random, are doors that take $6$ seconds to
cross instead of $1$.

\omcode{code/grade-12/05-graph-algorithms/listings/maze_doors.py}{1}{21}{The maze with doors.}%
\label{lst:g12:graph-algorithms:doors}

\omoutput{code/grade-12/05-graph-algorithms/outputs/maze_doors.txt}{The path with fewest steps is not the fastest.}%
\label{out:g12:graph-algorithms:doors}

\needspace{4\baselineskip}\medskip\noindent\textbf{Part I --- A maze as a graph.}
\begin{enumerate}
  \item What are the vertices and the edges of the maze's \omterm{def:g10:graphs-first-look:graph}{graph}? How many vertices?
  \item The carving makes a tree of passages. How many passages does it open? How many does the program count in the
    end, and why?
  \item In \code{make_maze}, which structure of the chapter does \code{stack} implement, and which search does the
    carving perform?
  \item Why does the maze have several routes from the entrance to the exit, and not just one?
\end{enumerate}

\needspace{4\baselineskip}\medskip\noindent\textbf{Part II --- Searching for the exit.}
\begin{enumerate}[resume]
  \item How many squares does the depth-first search visit before it reaches the exit? Why so many?
  \item Explain the order of the four turns in \code{wall_follower}.
  \item How many steps does the hand on the wall take? Why more than the shortest exit?
  \item Why is the hand on the wall sure to reach the exit here, although the maze has loops?
  \item Which information does the wall follower keep, and which does a search keep?
\end{enumerate}

\needspace{4\baselineskip}\medskip\noindent\textbf{Part III --- The shortest exit.}
\begin{enumerate}[resume]
  \item How long is the shortest exit? Check it on the drawing by counting the stars.
  \item Why is the path built from the parents a shortest one?
  \item Is it the only shortest exit? How would you find out?
\end{enumerate}

\needspace{4\baselineskip}\medskip\noindent\textbf{Part IV --- Doors.}
\begin{enumerate}[resume]
  \item How are the doors stored in the \omterm{def:g10:graphs-first-look:directed}{weighted graph} \code{g}?
  \item Why can breadth-first search not find the fastest exit?
  \item How long does the fewest-steps exit take, and how many doors does it cross?
  \item How long is the fastest exit, in steps and in seconds, and how many doors does it cross?
  \item Why is \omterm{met:g12:graph-algorithms:dijkstra}{Dijkstra's algorithm} correct here?
  \item If every door took $1$ second, what would the fastest exit cost?
  \item With $n$ squares, what is the cost of the scan version of \omterm{met:g12:graph-algorithms:dijkstra}{Dijkstra's algorithm}, and is it acceptable here?
  \item State the result: the shortest exit against the hand on the wall, and the fastest exit with doors.
\end{enumerate}
\end{problem}
